\chapter{Die fraktale Regularisierung der QED}

Die Quantenelektrodynamik (QED) hat divergente Integrale (Selbstenergie, Lamb-Shift), die durch Renormierung beseitigt werden. Die WDBT+ benötigt keine Renormierung – die fraktale Gewichtung wirkt als \emph{natürlicher Cutoff}.

\subsection{Die Selbstenergie des Elektrons}

Die Selbstenergie des Elektrons wird zu:

\begin{equation}
\Delta E_{\text{self}}^{\text{WDBT+}} = \frac{\alpha \hbar c}{\pi} \int_0^{\infty} \frac{dk}{k} \cdot \mathcal{F}(k) \cdot \left( \frac{k}{k_*} \right)^{D-3}.
\end{equation}

Für \( D \approx 2,704 \) ist \( D-3 \approx -0,296 \), und das Integral konvergiert. Der Cutoff ist \( k_* = 1/l_* \), wobei \( l_* \) die mesoskopische Skala ist.

\subsection{Der Lamb-Shift}

Der Lamb-Shift wird zu:

\begin{equation}
\Delta E_{\text{Lamb}}^{\text{WDBT+}} = \Delta E_{\text{Lamb}}^{\text{QED}} \cdot \left( \frac{a_0}{l_*} \right)^{D-3}.
\end{equation}

Die Kalibration mit dem experimentellen Lamb-Shift (\( \Delta E_{\text{Lamb}}^{\text{exp}} \approx 4,37 \times 10^{-6} \) eV) liefert \( l_* \approx a_0 \) – die atomare Skala.

\subsection{Der numerische Test}

Die numerische Integration des fraktal gewichteten Integrals reproduziert den gemessenen Lamb-Shift – ohne Renormierung, ohne willkürlichen Cutoff.